\chapter{R5：闭合历史记忆如何进入四象限}

\section{本章要解决的困惑}

很多文档把 R5 写得很神秘，好像它是订单簿里某种复杂微观结构。当前 runtime 里不是这样。
当前 R5 是一个\term{已闭合交易结果的短记忆}。它只看最近已经结束的 shadow entries，不看当前未来。

\warningline{R5 不是 entry 后未来 label，也不是当前 L2 depth 因子。它来自已经闭合的历史 outcome。}

\section{ClosedOutcomeState 的定义}

代码入口：

\begin{lstlisting}
systems/ccusdt_replay_exchange/strategies/python/ccusdt_tfi_core_idle01/online_features.py
ClosedOutcomeState
\end{lstlisting}

窗口长度是 5。设最近已经闭合的 5 笔收益为：

\[
g_1,g_2,g_3,g_4,g_5.
\]

代码计算：

\[
P_5=\sum_{i=1}^{5} \max(g_i,0),
\quad
N_5=\sum_{i=1}^{5} \max(-g_i,0).
\]

R5 是：

\[
R5=\frac{P_5}{N_5+\epsilon}.
\]

这里 \(g_i\) 是已经闭合的 PnL bps。不是未来，也不是当前 entry 的收益。

\section{手算例子}

假设最近 5 笔已经闭合结果是：

\begin{longtable}{p{0.18\textwidth}p{0.24\textwidth}p{0.46\textwidth}}
\toprule
index & pnl bps & 归类 \\
\midrule
1 & +3.0 & positive \\
2 & -1.0 & negative \\
3 & +2.0 & positive \\
4 & +1.5 & positive \\
5 & -0.5 & negative \\
\bottomrule
\end{longtable}

那么：

\[
P_5=3.0+2.0+1.5=6.5,
\quad
N_5=1.0+0.5=1.5.
\]

\[
R5=\frac{6.5}{1.5}=4.33.
\]

所以：

\[
A_t = 1\{R5_t \ge 1\}=1.
\]

如果最近 5 笔是：

\[
[-2.0,-1.5,+0.5,-1.0,+0.2],
\]

则：

\[
P_5=0.7,\quad N_5=4.5,\quad R5=0.156.
\]

此时 \(A_t=0\)。

\section{四象限的两个轴}

当前 cell 只看两个布尔量：

\[
A_t = 1\{R5_t \ge 1\},
\quad
B_t = 1\{frames\_since\_mid\_change_t \ge 31\}.
\]

于是：

\begin{longtable}{p{0.18\textwidth}p{0.18\textwidth}p{0.25\textwidth}p{0.29\textwidth}}
\toprule
\(A_t\) & \(B_t\) & cell & 人话 \\
\midrule
0 & 0 & \code{00_none} & R5 不好，mid 不 stale \\
1 & 0 & \code{10_r5_only} & R5 好，但 mid 不 stale \\
0 & 1 & \code{01_frames_only} & R5 不好，但 mid stale \\
1 & 1 & \code{11_r5_frames} & R5 好，且 mid stale \\
\bottomrule
\end{longtable}

这就是“四象限”。它不是四个复杂模型，只是两个 yes/no 切口。

\section{R5 在当前策略中的真实角色}

R5 不决定方向。方向已经由 TFI 决定：

\[
side_t = \operatorname{sign}(TFI_t).
\]

R5 也不是 entry 的第一开关。entry 候选已经由 TFI/past event/frames/membership 产生。

R5 的主要作用是：

\begin{lstlisting}
cell = _cell(a_r5, b_frames)
gamma = GAMMA_PRESETS[gamma_preset][cell]
target_exposure = weak_overlay_weight * gamma
\end{lstlisting}

换句话说，R5 在当前 runtime 里是 sizing/context，不是新的 entry primitive。

\section{Delta、Energy、Z 为什么容易混乱}

\code{ClosedOutcomeState} 里也有：

\begin{lstlisting}
delta  = positives - negatives
energy = positives + negatives
z      = (positive_bps + negative_bps) / sqrt(abs(positive_bps)+abs(negative_bps)+eps)
\end{lstlisting}

这些在研究上有意义，因为 R5 这个比例会丢掉强度信息。但必须诚实：

\warningline{\code{Delta_bps / Energy_bps} 目前没有进入 CC runtime 控制。}

当前 runtime cell 只看 \(R5 \ge 1\) 和 frames 阈值。不要读了研究报告就以为 Delta/Energy 已经在 Bot 里控制 sizing。

\section{手动检查点}

拿一组 closed outcomes：

\begin{lstlisting}
closed_bps = [2.0, -1.0, 0.5, 1.5, -0.2]
frames_since_mid_change = 34
\end{lstlisting}

手算：

\[
P_5=2.0+0.5+1.5=4.0,\quad
N_5=1.0+0.2=1.2,\quad
R5=3.33.
\]

\[
A=1,\quad B=1,\quad cell=\code{11_r5_frames}.
\]

这就是四象限计算的全部。

